\documentclass[conference]{IEEEtran}
\IEEEoverridecommandlockouts

% ── Encoding & font ──────────────────────────────────────────────────────────
\usepackage[utf8]{inputenc}
\usepackage[T1]{fontenc}

% ── Core IEEE packages ────────────────────────────────────────────────────────
\usepackage{cite}
\usepackage{amsmath,amssymb,amsfonts}
\usepackage{graphicx}
\usepackage{textcomp}
\usepackage{xcolor}

% ── Additional packages ───────────────────────────────────────────────────────
\usepackage{booktabs}          % professional table rules
\usepackage{caption}           % caption* for unnumbered captions

% ── Graphics path (relative to this file) ────────────────────────────────────
\graphicspath{{figures/}}

\def\BibTeX{{\rm B\kern-.05em{\sc i\kern-.025em b}\kern-.08em
    T\kern-.1667em\lower.7ex\hbox{E}\kern-.125emX}}

% ── Convenience macros ────────────────────────────────────────────────────────
\newcommand{\qreal}{Q_{\mathrm{real}}}
\newcommand{\qvirt}{Q_{\mathrm{virtual}}}
\newcommand{\pacb}{p_{\mathrm{acb}}}
\newcommand{\ema}{\mathrm{EMA}}
\newcommand{\lam}{\lambda}

% =============================================================================
\begin{document}
% =============================================================================

\title{A Proactive Digital Twin Architecture with Continuous Calibration\\for Mitigating 5G RRC Signaling Storms}

\author{%
\IEEEauthorblockN{1\textsuperscript{st} Ömer Dikyol}
\IEEEauthorblockA{\textit{Department of Computer Engineering}\\
\textit{[Istanbul Technical University]}\\
Istanbul, Turkey\\
dikyol25@itu.edu.tr}
}

\maketitle

% =============================================================================
% ABSTRACT
% =============================================================================
\begin{abstract}
In 5G New Radio architecture, the gNodeB (gNB) Random Access Channel (RACH) infrastructure has a structural vulnerability to high-volume RRC Signaling Storm traffic generated by coordinated botnet attacks. The current standard Access Class Barring (ACB) mechanism suffers from control loop latency due to its reactive nature that only activates after buffer saturation, proving insufficient under burst conditions. Predictive Digital Twin (DT) based approaches address this latency; however, the Model Drift phenomenon that emerges under non-stationary traffic conditions causes State Mirroring Error to accumulate over time, degrading prediction quality. This paper proposes a Proactive Digital Twin architecture with a continuous calibration cycle that actively suppresses Model Drift. To model the physical gNB RACH, a two-state Continuous-Time Markov Chain (CTMC) based MMPP/M/c/K queue is used; proactive control decisions are derived from an adaptive Exponential Moving Average (EMA) predictor. 1000-second Monte Carlo experiments conducted on N=10 independent random seeds confirm that the calibration cycle reduces State Mirroring Error by 42% compared to an uncalibrated DT. Connection success rate is statistically equivalent to the reactive baseline; P90 and P99 queueing delay values are reduced.
\end{abstract}

\begin{IEEEkeywords}
5G NR, RRC Signaling Storm, Digital Twin, Access Class Barring, MMPP/M/c/K,
Exponential Moving Average, Model Drift, State Mirroring Error, Queueing Delay
\end{IEEEkeywords}

% =============================================================================
% SECTION 1: INTRODUCTION
% =============================================================================
\section{Introduction}

The 5G New Radio (NR) architecture positions Radio Resource Control (RRC) as a central resource management layer to accommodate increasing user density and diversified service requirements. The gNodeB (gNB) processes RRC connection requests from user devices through the Random Access Channel (RACH) and responds to these requests through a specified number of preamble channels. Under normal traffic this design provides efficient resource allocation, but it leaves a security gap against coordinated botnet attacks.

Coordinated botnet traffic generates high-intensity, bursty RRC connection request waves in short intervals, intentionally saturating the RACH buffer. Such attacks are defined in the literature as \emph{RRC Signaling Storm}~\cite{b1} and lead to exhaustion of gNB-based resources, resulting in rejection of legitimate user connection requests. Considering the 5G NR's promise of low latency and high reliability for applications such as autonomous vehicles, remote surgery, and industrial automation, such service interruptions represent critical Quality of Service (QoS) failures.

The Access Class Barring (ACB) mechanism defined in 3GPP standards is the primary reactive tool for limiting RACH congestion. However, this mechanism has a fundamental \emph{Control Loop Latency} problem: the barring decision can only be made after buffer saturation is observed, thus losing its effectiveness precisely during the load increase period that needs to be prevented~\cite{b2}.

To overcome this limitation, predictive control approaches have attracted growing attention. The Digital Twin (DT) paradigm enables maintaining a virtual copy of the physical system in real-time and deriving control decisions from that copy~\cite{b3}. However, standard DT-based predictors suffer from \emph{Model Drift} under non-stationary burst conditions; the deviation between the state representations of the physical system and its virtual copy, the \emph{State Mirroring Error}, accumulates over time and degrades the effectiveness of the proactive controller.

This work addresses this gap by proposing a \emph{Calibrated Digital Twin} architecture with a continuous calibration cycle. The main contributions are:

\begin{enumerate}

\item \textbf{Realistic traffic model:} A two-state MMPP/M/c/K queue is designed to model 5G RRC Signaling Storm conditions; the burst dynamics of adversarial traffic are captured through Continuous-Time Markov Chain (CTMC).

\item \textbf{Continuous calibration cycle:} A calibration mechanism is presented that updates the Exponential Moving Average (EMA) weight using an asymmetric step rule when State Mirroring Error exceeds a threshold; this mechanism reduces State Mirroring Error by 42\%, as demonstrated through simulation.

\item \textbf{Proactive QoS improvement:} The effects of the calibration cycle on $P_{success}$ and Queueing Delay distribution (Cumulative Distribution Function) are measured through Monte Carlo experiments across $N=10$ independent seeds.

\item \textbf{Epsilon sensitivity analysis:} The impact of the calibration threshold $\varepsilon$ on system performance is systematically analyzed, demonstrating that the proposed $\varepsilon=10$ value is near-optimal.

\end{enumerate}

% =============================================================================
% SECTION 2: RELATED WORK
% =============================================================================
\section{Related Work}

\subsection{Reactive Methods and Limitations}

The most common reactive mechanism for preventing RACH congestion is the Access Class Barring (ACB) method defined in 3GPP TS 38.331~\cite{b4}. The original ACB framework defined in LTE Release~11 was based on a binary barring mechanism that grouped devices by access class (AC~0--9): the gNB would randomly delay RACH requests from certain access classes by publishing a barring factor $p$ and barring duration $T_{acb}$ through System Information Block 2 (SIB2).

With 5G NR Release~15, 3GPP adopted a more flexible framework called \emph{Unified Access Control} (UAC)~\cite{b4}. UAC presents a two-dimensional policy table that cross-matches traffic categories with access identities; published barring factors are announced to all devices through the \texttt{uac-BarringInfo} field of System Information Block 1 (SIB1). The $\pacb$ value calculated by the proposed DT controller at each $\Delta$ step directly corresponds to this UAC barring factor: the DT is designed to inject $\pacb$ into the gNB's SIB1 broadcast, providing a standard-compliant real-time interface. Thus, the proposed proactive control mechanism can be integrated into existing 3GPP signaling infrastructure without any protocol modifications.

This mechanism has been extensively examined in both LTE and 5G NR contexts; adaptive ACB extensions that adjust barring parameters based on queue occupancy ratio have been proposed~\cite{b5}.

However, all reactive ACB and UAC variants share a fundamental causality constraint: the barring decision depends on observing the queue length, processing this observation, and transmitting the control signal via SIB1. This process inevitably produces Control Loop Latency. Kotuliak et al.~\cite{b6}, in their measurement study on commercial 5G NR Non-Standalone networks, reported that the RACH procedure exceeded its latency budget under high load; this confirmed that reactive control loops are insufficient in practice.

\subsection{Proactive and Predictive Approaches}

In response to the inadequacy of reactive methods, predictive control approaches have been developed that apply barring policy before congestion occurs. Approaches in this area include Markov decision processes~\cite{b7}, reinforcement learning-based bandwidth allocation~\cite{b8}, and load predictors using time series regression~\cite{b9}.

The Digital Twin (DT) paradigm provides a broader framework that continuously updates the virtual representation of the physical system with real-time sensor feedback~\cite{b10}. DT applications have been studied in 5G network management for energy efficiency optimization~\cite{b11}, slice management~\cite{b12}, and anomaly detection~\cite{b13}. However, the vast majority of existing DT-based network controllers do not perform model updates during operation. This inevitably leads to Model Drift in botnet-driven RRC Signaling Storm scenarios, degrading prediction quality. The proposed \emph{Calibrated DT} architecture addresses this gap.

% =============================================================================
% SECTION 3: SYSTEM MODEL
% =============================================================================
\section{System Model}

\subsection{Network Environment and Attack Scenario}

This paper considers an RRC Signaling Storm scenario targeting a single gNB cell operating within the 5G NR standard. The attack consists of high-intensity, bursty RRC connection requests generated by a coordinated botnet, aiming to saturate the physical gNB buffer. Results are compared against a reactive baseline.

\subsection{Traffic Model: MMPP/M/c/K Queue}

To model the sudden transitions of adversarial traffic, a Markov-Modulated Poisson Process (MMPP) based on a two-state CTMC has been adopted. The system is characterized as $MMPP/M/c/K$: the arrival process is MMPP, service times are exponentially distributed, there are $c$ parallel servers, and the total system capacity is limited to $K$.

The CTMC defines two states. $S_0$ (Normal state): legitimate RRC connection requests are generated at $\lambda_{\mathrm{low}} = 10$ requests/s rate. $S_1$ (Burst state): adversarial botnet traffic is generated at $\lambda_{\mathrm{high}} = 500$ requests/s rate. Transitions between the two states are defined by the Generator Matrix $Q$:

\begin{equation}
Q = \begin{bmatrix}
    -\alpha_{mmpp} & \alpha_{mmpp} \\
    \beta_{mmpp}   & -\beta_{mmpp}
\end{bmatrix}
\label{eq:generator}
\end{equation}

\noindent where $\alpha_{mmpp} = 0.1\ \mathrm{s}^{-1}$ represents the $S_0 \to S_1$ transition rate, and $\beta_{mmpp} = 0.5\ \mathrm{s}^{-1}$ represents the $S_1 \to S_0$ transition rate. The State Sojourn Time spent in state $i$ is sampled from the $\mathrm{Exp}(|Q_{ii}|)$ distribution. This approach exactly reproduces CTMC dynamics in a discrete-event simulation environment without requiring matrix exponentiation calculations.

\subsection{Physical gNB Constraints}

The physical gNB RACH consists of $c = 54$ preamble channels. The total system capacity is limited to $K = 200$. The instantaneous queue length $\qreal$ is defined as the sum of requests in service and those waiting. When $\qreal(t) \geq K$ condition occurs, incoming requests are directly dropped due to buffer overflow.

% =============================================================================
% SECTION 4: PROPOSED METHOD
% =============================================================================
\section{Proposed Method}

\subsection{Proactive Digital Twin Architecture}

The proposed method consists of a Proactive Digital Twin (DT) controller operating in parallel with the physical gNB. The DT executes a combined pipeline with three consecutive stages at $\Delta = 1\ \mathrm{s}$ control interval: (i) observation and prediction, (ii) continuous calibration, (iii) control decision.

\subsection{EMA-based Arrival Rate Estimation}

The DT uses the Exponential Moving Average (EMA) filter to predict future arrival rates:

\begin{equation}
\ema_t = \alpha \cdot \lambda_{\mathrm{current}} + (1 - \alpha) \cdot \ema_{t-1}
\label{eq:ema}
\end{equation}

\noindent where $\lambda_{\mathrm{current}} = n_{\mathrm{window}}/\Delta$ represents the observed arrival count in the $\Delta$ window, and $\alpha \in [0,1]$ represents the dynamic correction weight. The initial value of $\alpha$ is set to $\alpha_0 = 0.2$, which falls within the $[0.15,\ 0.25]$ range recommended in the literature for bursty traffic~\cite{b8}.

\subsection{Stochastic Control Policy: $p_{acb}$ Calculation}

The DT analytically calculates the minimum dropping probability needed to prevent $\qreal$ from exceeding $K$ in the next $\Delta$ step:

\begin{equation}
\pacb = \max\!\left(0,\ \min\!\left(1,\ 1 - \frac{K - \qreal}
    {\max(\ema_t \cdot \Delta,\ \varepsilon_g)}\right)\right)
\label{eq:pacb}
\end{equation}

\noindent where $\varepsilon_g = 10^{-9}$ provides zero-division protection. $\pacb$ is applied as an independent Bernoulli trial to each incoming request, forming a random barring mechanism compatible with the 3GPP TS 38.331-defined Unified Access Control framework.

\subsection{Virtual Queue Update}

The DT maintains a virtual queue $\qvirt$ to internally simulate the controlled physical system:

\begin{equation}
\qvirt(t{+}\Delta) = \mathrm{clip}\!\left(
    \qvirt(t) + \ema_t \cdot (1 - \pacb) \cdot \Delta - c \cdot \Delta,\;
    0,\; K\right)
\label{eq:qvirtual}
\end{equation}

The presence of the $(1 - \pacb)$ factor is mathematically necessary. When the control policy is applied, the actual number of requests entering the system is not $\ema_t \cdot \Delta$ but $\ema_t \cdot (1 - \pacb) \cdot \Delta$. When this factor is excluded, $\qvirt$ permanently converges to $K$ during burst periods, artificially inflating State Mirroring Error and incorrectly triggering the calibration cycle.

\subsection{State Mirroring Error}

The instantaneous deviation between the DT and physical system is measured by State Mirroring Error $E$:

\begin{equation}
E(t) = \left| \qreal(t) - \qvirt(t) \right|
\label{eq:error}
\end{equation}

\subsection{Continuous Calibration Cycle}

After calculating $E$ at each $\Delta$ step, the DT updates the $\alpha$ weight using an asymmetric step rule:

\begin{equation}
\alpha \leftarrow \begin{cases}
    \min(\alpha_{max},\ \alpha + \delta^{+}) & \text{if } E > \varepsilon \\
    \max(\alpha_{min},\ \alpha - \delta^{-}) & \text{otherwise}
\end{cases}
\label{eq:calibration}
\end{equation}

\noindent In this work, $\varepsilon = 10$, $\delta^{+} = 0.05$, $\delta^{-} = 0.01$, $\alpha_{min} = 0.05$, and $\alpha_{max} = 1.0$ are determined. The intentional asymmetry in step sizes reflects two different operational requirements. The rapid $\delta^{+}$ increase triggered at attack onset transforms the EMA filter into a high-cut frequency structure, rapidly reducing the weight of past observations. When traffic normalizes, the slow decay with small $\delta^{-}$ maintains statistical stability.

\subsection{Combined Control Cycle}

All the above-defined components are integrated in the combined control cycle given in Algorithm~\ref{alg:cdt}. The cycle is triggered every $\Delta$ seconds, consuming only $\qreal$ and $n_{\mathrm{win}}$ inputs, and producing $\pacb$, updated $\qvirt$, and new $\alpha$ values as outputs.

\begin{figure}[t]
\vspace{-2pt}
\noindent\rule{\columnwidth}{0.5pt}\\[-5pt]
\noindent\rule{\columnwidth}{0.15pt}
\vspace{3pt}
{\small
\noindent\textbf{Algorithm 1:}
  Calibrated Digital Twin Control Cycle\\[4pt]
\noindent\textbf{Input:}
  $\qreal$,\ $n_{\mathrm{win}}$,\ $\Delta$,\ $\varepsilon$,\ $\alpha$,\
  $\alpha_{min}$,\ $\alpha_{max}$,\ $\delta^{+}$,\ $\delta^{-}$\\
\noindent\textbf{Output:}
  $\pacb$,\ updated $\qvirt$,\ current $\alpha$\\[4pt]
\noindent\textbf{At each} $\Delta$ step:\\[2pt]
\noindent\phantom{0}1:\enspace
  $\lambda_{cur} \leftarrow n_{\mathrm{win}} / \Delta$
  \hfill\textit{\% instantaneous arrival rate}\\
\noindent\phantom{0}2:\enspace
  $\ema_t \leftarrow \alpha\cdot\lambda_{cur}+(1-\alpha)\cdot\ema_{t-1}$
  \hfill\textit{\% EMA update}\\
\noindent\phantom{0}3:\enspace
  $d \leftarrow \max\!\left(\ema_t\cdot\Delta,\ 10^{-9}\right)$\\
\noindent\phantom{0}4:\enspace
  $\pacb \leftarrow \max\!\left(0,\,\min\!\left(1,\,1-(K-\qreal)/d\right)\right)$
  \hfill\textit{\% \eqref{eq:pacb}}\\
\noindent\phantom{0}5:\enspace
  $\qvirt \leftarrow
   \mathrm{clip}\!\left(\qvirt+\ema_t(1{-}\pacb)\Delta-c\Delta,\;0,\;K\right)$
  \hfill\textit{\% \eqref{eq:qvirtual}}\\
\noindent\phantom{0}6:\enspace
  $E \leftarrow \left|\qreal - \qvirt\right|$
  \hfill\textit{\% \eqref{eq:error}}\\
\noindent\phantom{0}7:\enspace
  \textbf{if} $E > \varepsilon$ \textbf{then}\\
\noindent\phantom{07:}\enspace
  $\alpha \leftarrow \min(\alpha_{max},\;\alpha + \delta^{+})$
  \hfill\textit{\% fast adaptation}\\
\noindent\phantom{0}8:\enspace
  \textbf{else}\\
\noindent\phantom{08:}\enspace
  $\alpha \leftarrow \max(\alpha_{min},\;\alpha - \delta^{-})$
  \hfill\textit{\% slow decay}\\
\noindent\phantom{0}9:\enspace
  \textbf{end if}
  \hfill\textit{\% \eqref{eq:calibration}}\\
\noindent10:\enspace
  $\text{drop\_rate} \leftarrow \pacb$
  \hfill\textit{\% transfer control decision to gNB}
}
\vspace{3pt}
\noindent\rule{\columnwidth}{0.15pt}\\[-5pt]
\noindent\rule{\columnwidth}{0.5pt}
\label{alg:cdt}
\end{figure}

% =============================================================================
% SECTION 5: PERFORMANCE EVALUATION
% =============================================================================
\section{Performance Evaluation}

\subsection{Simulation Parameters}

Performance evaluation is conducted through Monte Carlo experiments run for $T_{\mathrm{sim}} = 1000\ \mathrm{s}$ duration on $N = 10$ independent random seeds. The simulation was implemented using Python 3.10 with the SimPy library via discrete-event method; results were analyzed with NumPy and Pandas. Key parameters are summarized in Table~\ref{tab:params}.

\begin{table}[t]
\caption{Simulation Parameters}
\label{tab:params}
\begin{center}
\begin{tabular}{lll}
\toprule
\textbf{Parameter} & \textbf{Symbol} & \textbf{Value} \\
\midrule
Total duration              & $T_{\mathrm{sim}}$ & 1000 s \\
RACH preamble channels      & $c$                & 54 \\
Buffer capacity             & $K$                & 200 \\
Normal arrival rate         & $\lambda_{low}$    & 10 requests/s \\
Burst arrival rate          & $\lambda_{high}$   & 500 requests/s \\
Normal$\to$Burst transition & $\alpha_{mmpp}$    & 0.1 s$^{-1}$ \\
Burst$\to$Normal transition & $\beta_{mmpp}$     & 0.5 s$^{-1}$ \\
Control interval            & $\Delta$           & 1 s \\
Calibration threshold       & $\varepsilon$      & 10 \\
Initial EMA weight          & $\alpha_0$         & 0.20 \\
Reactive threshold          & $0.8K$             & 160 \\
Reactive drop rate          & ---                & 0.50 \\
\bottomrule
\end{tabular}
\end{center}
\end{table}

\subsection{State Mirroring Error Analysis}

Table~\ref{tab:metrics} summarizes the key performance metrics of the three methods across $N=10$ seeds. Confidence intervals are calculated using the formula $\bar{x} \pm 1.96 \cdot \hat{\sigma}/\sqrt{N}$.

\begin{table}[t]
\caption{Multi-seed Performance Metrics ($N=10$,\ $T_{\mathrm{sim}}=1000$~s)}
\label{tab:metrics}
\begin{center}
\begin{tabular}{lccc}
\toprule
\textbf{Method}
    & $P_{success}$       & $P_{drop}$          & $E_{mean}$         \\
\textbf{}
    & ($\pm$ 95\% CI)     & ($\pm$ 95\% CI)     & ($\pm$ 95\% CI)    \\
\midrule
Reactive Baseline       & $0.310 \pm 0.015$ & $0.688 \pm 0.015$ & $0.00 \pm 0.00$ \\
DT (No Calibration)     & $0.301 \pm 0.011$ & $0.698 \pm 0.011$ & $85.03 \pm 4.40$ \\
Calibrated DT           & $0.309 \pm 0.019$ & $0.690 \pm 0.019$ & $49.14 \pm 3.09$ \\
\bottomrule
\end{tabular}
\end{center}
\end{table}

The State Mirroring Error results confirm the calibration benefit. The uncalibrated DT has $E_{mean} = 85.03$. With the continuous calibration cycle enabled, $E_{mean}$ drops to $49.14$, a \textbf{42\% reduction}. The $E_{mean}$ distributions of the two DT configurations do not overlap, confirming a statistically significant difference. Figure~\ref{fig:E_ci} presents the State Mirroring Error time series with 95\% Confidence Interval bands.

\begin{figure}[t]
\centerline{\includegraphics[width=\columnwidth]{phase5_E_ci.png}}
\caption{State Mirroring Error time series: comparison between Calibrated DT (blue) and DT (No Calibration) (yellow). Shaded regions show 95\% Confidence Interval across $N=10$ seeds. The horizontal dashed line represents the calibration threshold $\varepsilon=10$.}
\label{fig:E_ci}
\end{figure}

\subsection{Connection Success Rate and Queue Stability}

Regarding connection success rates, the following ranking is observed: Reactive Baseline ($P_{success} = 0.310$), Calibrated DT ($P_{success} = 0.309$), and DT (No Calibration) ($P_{success} = 0.301$). The difference in $P_{success}$ between the Reactive Baseline and Calibrated DT ($\Delta = 0.001$) is statistically insignificant due to wide overlap of confidence intervals. However, this numerical similarity should not be read as mechanistic equivalence.

The Reactive Baseline achieves a similar $P_{success}$ value by activating the dropping policy when buffer saturation approaches (i.e., after the $\qreal > 0.8K$ threshold is exceeded). In contrast, the Calibrated DT gradually increases $\pacb$ to keep $\qreal$ in a stable region before it reaches $K$. Figure~\ref{fig:q_real_ci} presents the Queue Stability comparison.

\begin{figure}[t]
\centerline{\includegraphics[width=\columnwidth]{phase5_q_real_ci.png}}
\caption{Queue length $\qreal$ time series: comparison of three methods with 95\% Confidence Interval bands. The upper dashed line shows the $K=200$ buffer limit; the lower dashed line shows the $0.8K=160$ reactive threshold.}
\label{fig:q_real_ci}
\end{figure}

\subsection{Queueing Delay: Cumulative Distribution Function Analysis}

The per-request Queueing Delay distribution shows the advantage of the Calibrated DT more directly. The Cumulative Distribution Function (CDF) in Figure~\ref{fig:delay_cdf} pools delay values from all served requests across $N=10$ seeds.

\begin{figure}[t]
\centerline{\includegraphics[width=\columnwidth]{phase5_delay_cdf.png}}
\caption{Per-request Queueing Delay Cumulative Distribution Function (CDF): all seeds pooled. Solid dots show P50, P90, and P99 percentiles. The leftward shift of the Calibrated DT curve indicates a lower delay distribution.}
\label{fig:delay_cdf}
\end{figure}

The CDF curve for the Calibrated DT shifts left relative to both baselines. As buffer saturation approaches, the Reactive Baseline causes requests to accumulate in the waiting queue, causing sharp jumps in P90 and P99 delay values. The Calibrated DT keeps $\qreal$ in a stable region, reducing the delay tail for served requests.

\subsection{Epsilon Sensitivity Analysis}

Sensitivity analysis over $\varepsilon \in \{5, 10, 20, 40, 80\}$ shows that $\varepsilon = 10$ is near-optimal for both $P_{success}$ and the normalized objective score, confirming that the parameter selection is empirically grounded rather than arbitrary.

\subsection{Computational Complexity Analysis}

The computational load of the proposed Calibrated DT controller is a critical design constraint for its integration into actual gNB hardware. Examining Algorithm~\ref{alg:cdt}, each $\Delta$ step consists of a fixed number of basic arithmetic operations: one division (line~1), one multiplication-addition pair (line~2), one clipping operation, and a few comparisons. The stored state consists only of $\{\ema_t, \qvirt, \alpha, n_{\mathrm{win}}\}$.

The controller therefore requires $O(1)$ time and $O(1)$ space per control step. Fixed memory makes it compatible with the limited resources of gNB hardware.

This contrasts with learning-based controllers. In Machine Learning (ML) or Deep Neural Network (DNN) based approaches, inference cost scales with the weight matrix size to $O(W)$, where $W$ is the number of trained weights. LSTM and Transformer-based traffic prediction models commonly reach $W \sim 10^5$--$10^7$ parameters; at that scale, inference latency reaches the millisecond range and can be non-deterministic~\cite{b9}. gNB RACH controllers require responses in the microsecond range, making DNN-based controllers impractical for direct hardware deployment.

The proposed $O(1)$ Calibrated DT addresses this constraint by providing accurate online adaptation within a fixed hardware budget. Each $\Delta = 1\ \mathrm{s}$ step completes in nanoseconds on modern gNB processors, making it compatible with commercial gNB software stacks.

\subsection{O-RAN Integration and Real-time Deployment}

The Calibrated DT controller maps directly onto the Open Radio Access Network (O-RAN) architecture. In the layered reference architecture defined by the O-RAN Alliance, the \emph{Near-Real-Time RAN Intelligent Controller} (Near-RT RIC) manages control loops between 10~ms and 1~s~\cite{b14}. The $\Delta = 1\ \mathrm{s}$ control interval aligns with the upper end of that range; the $O(1)$ complexity shown in Algorithm~\ref{alg:cdt} indicates the controller can operate within the 10~ms lower bound as well.

In the deployment architecture, the Calibrated DT controller is implemented as an \emph{xApp} running on Near-RT RIC. The input channel for the xApp is the \textbf{E2} interface: the gNB reports cell-level RACH metrics ($\qreal$, preamble collision count, and arrival count per window $n_{\mathrm{win}}$) within the E2 Service Model for Key Performance Measurements (E2SM-KPM) framework. The xApp consumes these inputs to execute Algorithm~\ref{alg:cdt} and produces the calculated $\pacb$ value as output.

The produced $\pacb$ is transmitted to the gNB CU-CP through the \textbf{E2SM-RC} (RAN Control) interface. The CU-CP injects this value into the \texttt{uac-BarringInfo} field of the SIB1 broadcast, in full compatibility with the 5G NR Unified Access Control framework described in Section~II. Thus, the control loop closes entirely over the standard O-RAN and 3GPP signaling infrastructure, requiring no protocol modifications or special hardware additions.

This deployment model has two practical advantages. First, operator parameters (e.g., attack sensitivity level or the $\varepsilon$ threshold) received from the Non-RT RIC via the A1 interface can be updated at xApp runtime, allowing the same xApp to serve different network configurations. Second, the xApp abstraction decouples the controller from gNB vendor hardware, so the same application runs across vendors.

% =============================================================================
% SECTION 6: DISCUSSION AND METHOD LIMITATIONS
% =============================================================================
\section{Discussion and Method Limitations}

The results show the effectiveness of the proposed Calibrated DT architecture in the single-cell attack scenario modeled by the MMPP/M/c/K queue. However, as with any simulation study, the validity of the model relies on several simplifying assumptions. This section states those limitations explicitly and identifies where the architecture may not hold.

\textbf{Single-Cell Scope.} The used $MMPP/M/c/K$ model considers a single gNB cell independently. In real 5G deployments, handover signaling, UE mobility correlations, and inter-cell interference exist between neighboring cells. These factors complicate the $\qreal$ observation process in a way that cannot be fully represented by independent Markov chains. The 42\% reduction seen in the single-cell case may therefore not transfer directly to multi-cell deployments. Overcoming this limitation requires either a distributed architecture running separate DT instances for each cell or one that delegates inter-cell state sharing to a central xApp coordinator.

\textbf{Fixed Calibration Steps.} The proposed asymmetric calibration rule relies on fixed step sizes of $\delta^{+} = 0.05$ and $\delta^{-} = 0.01$. These values were determined based on intuitive reasoning and practical limits supported by the $\varepsilon = 10$ sensitivity analysis. However, fixed $\delta^{+}$ and $\delta^{-}$ create a security gap: slowly ramping attacks that keep botnet traffic just below $\varepsilon$ can bring $\qreal$ close to $K$ without triggering calibration. This is a gap against more deliberate adversaries. Replacing the fixed steps with a Reinforcement Learning (RL) agent that treats instantaneous State Mirroring Error as a reward signal is a direct path to addressing this limitation.

% =============================================================================
% SECTION 7: CONCLUSION
% =============================================================================
\section{Conclusion}

This paper proposed and evaluated a Calibrated Digital Twin architecture for proactively mitigating botnet-driven RRC Signaling Storm attacks against the 5G NR gNB RACH infrastructure. The core motivation is that current reactive ACB mechanisms are insufficient under burst conditions due to Control Loop Latency, while standard predictive DT approaches suffer from Model Drift under non-stationary traffic transitions, causing State Mirroring Error to accumulate over time.

The system measures State Mirroring Error $E = |\qreal - \qvirt|$ at each $\Delta = 1\ \mathrm{s}$ step and updates the EMA weight using the asymmetric calibration rule in \eqref{eq:calibration}. Results across $N=10$ independent seeds show that the calibration cycle reduces State Mirroring Error by 42\%, suppresses the tail of the queueing delay distribution, and holds connection success rate at a level statistically equivalent to the reactive baseline.

The results point to a design requirement: in non-stationary 5G environments, predictive control effectiveness depends on how tightly the virtual model tracks the physical system, not just on the control policy. Active suppression of State Mirroring Error is therefore a necessary condition for deploying DT-based proactive controllers in practice.

Future work will focus on two directions. First, extending the architecture to multi-cell environments. Second, learning the $\delta^{+}$ and $\delta^{-}$ step sizes dynamically using a Reinforcement Learning (RL) agent that uses State Mirroring Error as a reward signal.

% =============================================================================
% REFERENCES
% =============================================================================
\begin{thebibliography}{00}

\bibitem{b1}
D. Rupprecht, K. Kohls, T. Holz, and C. Pöpper,
``Breaking LTE on Layer Two,''
\emph{Proc. IEEE Symp. Security and Privacy (S\&P)},
San Francisco, CA, USA, May 2019, pp.~1121--1136.

\bibitem{b2}
B. H. Lee, H. S. Lee, S. Moon, and J. W. Lee,
``Enhanced Random Access for Massive-Machine-Type Communications,''
\emph{IEEE Internet of Things J.},
vol.~8, no.~8, pp.~7046--7064, Apr. 2021.

\bibitem{b3}
M. Grieves and J. Vickers,
``Digital Twin: Mitigating Unpredictable, Undesirable Emergent Behavior in
Complex Systems,''
\emph{Transdisciplinary Perspectives on Complex Systems},
F.-J. Kahlen, S. Flumerfelt, and A. Alves, Eds.
Cham: Springer, 2017, pp.~85--113.

\bibitem{b4}
3rd Generation Partnership Project (3GPP),
``NR; Radio Resource Control (RRC) Protocol Specification,''
TS~38.331, Rel.~17.3.0, Sep. 2022.

\bibitem{b5}
L. Miuccio, D. Panno, and S. Riolo,
``Joint Control of Random Access and Dynamic Uplink Resource Dimensioning
for Massive MTC in 5G NR Based on SCMA,''
\emph{IEEE Internet of Things J.},
vol.~7, no.~6, pp.~5042--5063, Jun. 2020.

\bibitem{b6}
M. Kotuliak, T. K. Todorov, A. S. A. Gamel, and P. Počta,
``Performance Evaluation of 5G New Radio Non-Standalone Network in Commercial Deployment,''
\emph{IEEE Access},
vol.~9, pp.~110689--110698, Aug. 2021.

\bibitem{b7}
S. K. Sharma and X. Wang,
``Toward Massive Machine Type Communications in Ultra-Dense Cellular IoT
Networks: Current Issues and Machine Learning-Assisted Solutions,''
\emph{IEEE Commun. Surveys Tuts.},
vol.~22, no.~1, pp.~426--471, 1st Qtr. 2020.

\bibitem{b8}
R. Dong, C. She, W. Hardjawana, Y. Li, and B. Vucetic,
``Deep Learning for Radio Resource Allocation With Diverse Quality-of-Service
Requirements in 5G,''
\emph{IEEE Trans. Wireless Commun.},
vol.~20, no.~4, pp.~2309--2324, Apr. 2021.

\bibitem{b9}
C. Zhang, P. Patras, and H. Haddadi,
``Deep Learning in Mobile and Wireless Networking: A Survey,''
\emph{IEEE Commun. Surveys Tuts.},
vol.~21, no.~3, pp.~2224--2287, 3rd Qtr. 2019.

\bibitem{b10}
H. X. Nguyen, R. Trestian, D. To, and M. Tatipamula,
``Digital Twin for 5G and Beyond,''
\emph{IEEE Commun. Mag.},
vol.~59, no.~2, pp.~10--15, Feb. 2021.

\bibitem{b11}
M. Almasan, M. Gürsoy, S. Troia, G. Maier, and A. Pattavina,
``Digital Twin Network: Opportunities and Challenges,''
\emph{Proc. IFIP/IEEE Int. Symp. Integrated Network Management (IM)},
Bordeaux, France, May 2021, pp.~727--731.

\bibitem{b12}
A. Masood, D. Lakew, and S. Cho,
``Security Challenges and Solutions for 5G Network Slicing,''
\emph{IEEE Commun. Standards Mag.},
vol.~4, no.~1, pp.~63--69, Mar. 2020.

\bibitem{b13}
R. Fontana and G. Carra,
``Digital Twin-Based Anomaly Detection in Cyber-Physical Systems,''
\emph{Proc. IEEE Int. Conf. Cyber-Physical Systems, Networks, and
Applications (CPSNA)},
Naples, Italy, Sep. 2021, pp.~1--6.

\bibitem{b14}
M. Polese, L. Bonati, S. D'Oro, S. Basagni, and T. Melodia,
``Understanding O-RAN: Architecture, Interfaces, Algorithms, Security,
and Research Challenges,''
\emph{IEEE Commun. Surveys Tuts.},
vol.~25, no.~2, pp.~1376--1411, 2nd Qtr. 2023.

\end{thebibliography}

\end{document}